\documentclass[border=10pt]{standalone}
\usepackage{latentfig}
\input{build/pattern-values}
\begin{document}
\begin{figurecanvas}
\begin{tikzpicture}[x=0.3cm, y=0.3cm]
\input{build/pattern.tikz}
% Which part of the sequence each band of rows and columns is.
\pgfmathsetmacro{\outstart}{\patternn+1}
\node[note, anchor=south] at (\patternn/2,\patternsize+0.2) {input list (keys)};
\node[note, anchor=south] at ({(\outstart+\patternsize)/2},\patternsize+0.2) {sorted output};
\node[note, anchor=south, rotate=90] at (-1.9,{\patternsize-\patternn/2}) {input list (queries)};
\node[note, anchor=south, rotate=90] at (-1.9,{(\patternsize-\outstart)/2}) {sorted output};
\node[head, anchor=base west] at (0,\patternsize+2.0) {Head \patternlayer.\patternhead, the critical head of seed 0, on one test list.};
\node[note, anchor=north west, align=left] at (\patternsize+1.2,\patternsize) {
  Each row is a query: the position\\ attending. Each column is a key:\\ the position attended to.\\[4pt]
  Shading: the attention weight,\\ darker for more.\\[4pt]
  {\color{spot}Red rings}: for each output, the\\ input position holding the value\\ it has just written.\\[4pt]
  Back-hatched: later positions,\\ hidden by the causal mask.\\[4pt]
  \textsc{s}: the separator.};
\end{tikzpicture}
\end{figurecanvas}
\end{document}
